\documentclass[12pt,a4paper,oneside]{article}
\usepackage{amsmath}
\newcommand{\sol}{{\it Resposta: }}
\newcommand{\punt}{{\it Puntuació: }}
\begin{document}
\font\capsal=cmbx12
{\capsal
\noindent{\centre}\hfill 1r de Batxillerat

\noindent{Matemàtiques}\hfill Curs acadèmic: \curs

\noindent{Nom i cognoms:}\hfill Professor: \docent{}
\vskip .1 cm\hrule}
\section{Nombres radicals}

\begin{enumerate}
    \item Resoleu les següents potències combinades

\begin{enumerate}[label=\alph*)]
    \item $\left(2^3\right)^{-2}=$
    \hfill \sol{$\dfrac{1}{64}$}
    
    \item $\left[\left(3^{-2}\right)^3\right]^{-1}=$
    \hfill \sol{$3^6$}
    \item $\dfrac{\left(2^{-2}\cdot3^2\right)^{-1}}
    {\left(6^{-1}\right)^2}=$
    \hfill \sol{$2^4$}
    \end{enumerate}

\item Feu-ne algunes més. Simplifiqueu fins a obtenir una única potència o un nombre.

    \begin{enumerate}[label=\alph*)]
    \item $\dfrac{\left(2^{-2}\cdot5^2\right)^2}
    {10^{-2}}=$
    \hfill \sol{$\dfrac{15625}{4}$}
    \item $\left(\dfrac{a^{-2}b^3}{a^3b^{-2}}\right)^{-1}=$
    \hfill \sol{$\dfrac{a^5}{b^5}$}
    \item $\dfrac{\left[(3^2)^{-2}\cdot(3^{-1})^3\right]^{-1}}
    {3^{-5}}=$     \hfill \sol{$\dfrac{$3^{12}$}
    %\item $\left[\dfrac{(x^{-2}y^3)^2}{(x^3y^{-1})^{-1}}\right]^{-1}=$ \hfill \sol{$=\dfrac{x}{y^5}$}


\end{document}
